\begin{helper}
Fix \(\varepsilon_2>0\).  Let \(\beta=1/2\), let
\(m:\mathbb N\to\mathbb N\) satisfy
\[
m(n)=M(n)=\noderef{TwoBiteNaturalReducedVertexCount}(n)
\]
for every \(n\), and put
\[
p(n)=\noderef{TwoBiteEdgeProbability}(\beta,n).
\]
Under the two-bite law
\(\noderef{TwoBiteEventProbability}(n,m(n),p(n))\), the event
\[
\begin{aligned}
&\forall r\in \operatorname{Fin}(m(n)),\quad
\noderef{WithinRelativeError}\!\left(
  \varepsilon_2,\ p(n)m(n),\
  \noderef{GraphDegreeCount}(\noderef{TwoBiteRedGraph}(\omega),r)
\right),\\
&\forall b\in \operatorname{Fin}(m(n)),\quad
\noderef{WithinRelativeError}\!\left(
  \varepsilon_2,\ p(n)m(n),\
  \noderef{GraphDegreeCount}(\noderef{TwoBiteBlueGraph}(\omega),b)
\right)
\end{aligned}
\]
holds with high probability in the sense of \(\noderef{WithHighProbability}\).
\end{helper}
\begin{proof}
Write \(L=\log n\), \(M=M(n)=m(n)\), and
\[
p=p(n)=\beta\sqrt{\frac{L}{n}}=\frac12\sqrt{\frac{L}{n}}.
\]
By \(\noderef{TwoBiteNaturalReducedVertexCount}\), together with
\(\noderef{TwoBiteReducedVertexCount}\) and \(\noderef{TwoBiteStretch}\), for
all sufficiently large \(n\),
\[
M=\left\lfloor\frac{n}{L^2}\right\rfloor,\qquad
\frac12\frac{n}{L^2}\le M\le \frac{n}{L^2}.
\tag{1}
\]
Hence \(M\to\infty\), and
\[
pM \ge \frac14\,\frac{\sqrt n}{L^{3/2}}\longrightarrow\infty,
\qquad
\log M=O(L).
\tag{2}
\]
In particular \(pM/\log M\to\infty\).

The two-bite sample weight factors as the red graph weight, the blue graph
weight, and the uniform injection weight by \(\noderef{TwoBiteSampleWeight}\).
The red and blue graph weights are the independent \(G(M,p)\) product
weights encoded by \(\noderef{GnpGraphWeight}\), and their total masses are
\(1\) by \(\noderef{GnpGraphWeightSumOne}\).  Thus, after summing over the
irrelevant blue graph and injection coordinates, the marginal law of
\(\noderef{TwoBiteRedGraph}(\omega)\) is exactly the product Bernoulli-\(p\)
law on unordered red pairs.  The same argument gives the identical marginal
law for \(\noderef{TwoBiteBlueGraph}(\omega)\).  The eventual normalization
and nonempty injection support for this rounded value of \(M\) are supplied
by \(\noderef{TwoBiteNaturalTotalMassOneEventually}\); changing finitely many
initial \(n\)'s does not affect \(\noderef{WithHighProbability}\).

Fix a red vertex \(r\in\operatorname{Fin}(M)\).  Since simple graphs have no
loops, the degree
\[
D_r=\noderef{GraphDegreeCount}(\noderef{TwoBiteRedGraph}(\omega),r)
\]
is a sum of the \(M-1\) independent Bernoulli-\(p\) edge indicators joining
\(r\) to the other red vertices.  Its mean is
\[
\mu=(M-1)p=pM-p.
\tag{3}
\]
For a binomial sum \(Y\) with mean \(\mu\), the elementary Chernoff
calculation gives, for every fixed \(0<\delta\le 1/2\),
\[
\Pr(|Y-\mu|>\delta\mu)\le 2e^{-c_\delta\mu}
\tag{4}
\]
with \(c_\delta>0\).  Indeed,
\(\mathbb E e^{\lambda Y}=(1-p+pe^\lambda)^{M-1}
\le \exp(\mu(e^\lambda-1))\).  Markov's inequality with
\(\lambda=\log(1+\delta)\) gives the upper tail
\(\exp(-\mu((1+\delta)\log(1+\delta)-\delta))\), and the same argument with
\(-\lambda\), \(\lambda=-\log(1-\delta)\), gives the lower tail
\(\exp(-\mu((1-\delta)\log(1-\delta)+\delta))\).  Taking the smaller of the
two positive exponents proves (4).

Let \(\delta=\min\{\varepsilon_2/4,1/2\}\).  Since \(M\to\infty\), we have
\(p\le(\varepsilon_2/4)pM\) for all sufficiently large \(n\).  If
\(|D_r-\mu|\le\delta\mu\), then using \(\mu\le pM\) and (3),
\[
|D_r-pM|
\le |D_r-\mu|+|\mu-pM|
\le \delta\mu+p
\le \frac{\varepsilon_2}{4}pM+\frac{\varepsilon_2}{4}pM
\le \varepsilon_2 pM.
\]
Thus failure of
\(\noderef{WithinRelativeError}(\varepsilon_2,pM,D_r)\) implies the
Chernoff tail in (4).  By (2) and (3), for all large \(n\) the mean \(\mu\)
is at least \((1/2)pM\), so
\[
\Pr\!\left(
\neg\noderef{WithinRelativeError}(\varepsilon_2,pM,D_r)
\right)
\le 2\exp(-c pM)
\tag{5}
\]
for a constant \(c=c(\varepsilon_2)>0\).

The same estimate holds for each blue vertex \(b\), with
\[
D_b=\noderef{GraphDegreeCount}(\noderef{TwoBiteBlueGraph}(\omega),b).
\]
There are \(M\) red vertices and \(M\) blue vertices.  A union bound and
(5) give total failure probability at most
\[
4M\exp(-c pM).
\tag{6}
\]
By (2), \(pM/\log M\to\infty\), so the logarithm of the right side is
\[
\log 4+\log M-cpM\longrightarrow -\infty.
\]
Consequently (6) tends to \(0\).  Equivalently, the probability of the
simultaneous red and blue base-degree concentration event tends to \(1\),
which is exactly the displayed \(\noderef{WithHighProbability}\) statement.
\end{proof}
